\chapter{Spannung des Szintillationsdetektors}

\section{Myonenrate}\label{sec:myonenratios}

Es wurde die Spannung des Szintillationsdetektors $U_\text{Szint}$ so eingestellt, dass möglichst alle einfallenden Myonen als Signal detektiert wurden, aber möglichst kein Rauschen. Um dies zu erreichen, wurde die Rate an einfallenden Myonen, die auf die Detektorfläche in einer Minute einfallen ausgerechnet, ausgehend von der Rate von Myonen aus kosmischer Strahlung von ca. $1$ Muon pro $\text{cm}^2$ und pro Minute \cite[145]{astroparticles}. Aus den Maßen $\SI{28.5(5)}{\centi\m^2}\times\SI{5.0(5)}{\centi\m^2}$ des rechteckigen Detektors ergibt sich als Detektorfläche $\mathcal{A}=\SI{1.43(14)e-4}{\m^2}$ eine Rate von $A_\text{Lit}=\SI{2.38(24)}{\per\s}$.

Um die Rate der gemessenen Teilcheneinfälle (es wird hierzu angenommen, dass alle einfallenden Teilchen, die gemessen werden, Myonen sind \cite[2]{515versuchsanleitung}) zu bestimmen, wurde mithilfe des Programms \texttt{fpexperiment} \cite[5-6]{515versuchsanleitung} und einer vorher bestimmten Anzahl an zu messenden Ereignissen $n_\mu$ für verschiedene Versorgungsspannungen $U_\text{Szint}$ die Zeit $t_\mu$ gemessen, bis diese Anzahl an Teilcheneinfällen vom Detektor gemessen wurden.
Im Falle der optimalen Sensitivität des Detektors für Myonen sollte diese so bestimmte Rate $a_\mu$ der theoretischen Myoneneinfallrate auf die Detektorfläche entsprechen.

Als Ablesefehler für den Wert von $U_\text{Szint}$ wurde wieder $1\%$ des anzeigten Wertes angenommen \cite[2]{515versuchsanleitung}, für die gemessenen Zeiten ein Fehler von $3\%$ abgeschätzt, da nicht klar ist, wie genau die Messung des zur Zeitmessung verwendeten unix-Standardprogramms \texttt{time} \cite{timeunix}(misst die Zeitdauer zwischen Start und Ende eines Programms, hier jeweils von \texttt{fpexperiment}) ist und auch \texttt{fpexperiment} eine sehr kurze Startzeit hatte bei jedem Aufrufen. Da für jedes der detektierten Ereignisse ein Punkt in der Konsole ausgegeben wird \cite[6]{515versuchsanleitung}, was auch so beobachtet wurde, kann davon ausgegangen werden, dass die Ereignisanzahl fehlerlos ist.

Die so gemessenen Daten sind in Tabelle \ref{tab:ratios_data} aufgeführt, die daraus berechneten Raten in Tabelle \ref{tab:ratios}. Als am besten zur erwarteten Myonenrate passende Spannung $U_\text{Szint}$ ergab sich $U_\text{Szint}=\SI{-1.550(16)}{\kilo\V}$. Für die Langzeitmessung wurde die Versorgungsspannung des Szintillationsdetektors auf diesen Wert eingestellt.

\section{Elektronenrate}\label{sec:electron_rate}

Um im Folgenden die Betriebsparameter Diskriminatorschwelle, $U_D$ sowie Verzögerung des Aufbaus einstellen zu können, wurde auch für einen Einfall von Elektronen die passende Szintillationsdetektor-Spannung eingestellt. Um diese zu bestimmen, wurde die $\ce{^90Sr}$-Quelle erneut in der Mitte des Holzspalts auf der Driftkammer und der Szintillationsdetektor orthogonal zu den Anodendrähten (d.h. parallel zum Holzschlitz) unterhalb der Driftkammer platziert (Position Szintillationsdetektor siehe Abbildung \ref{fig:scinti_downlow}).
Es wurde mithilfe von \texttt{fpexperiment} die Rate von gemessenen Teilcheneinfällen im Detektor für verschiedene $U_\text{Szint}$ bestimmt. Die Spannungseinstellung ist dann ideal, wenn das Verhältnis $A$ der Rate der detektierten Teilcheneinfälle $a_e=\frac{n_e}{t_e}$ für eine gegebene Spannungseinstellung mit $\ce{^90Sr}$-Quelle (hier werden als einfallende Teilchen Elektronen und Myonen gemessen) zur Rate der detektierten Teilcheneinfällen zur gleichen Spannungseinstellung ohne Quelle über dem Holzschlitz (hier werden nur die kosmischen Myonen als Untergrundereignisse gemessen) maximal wird. In diesem Fall sind die Ereignisse, die vom Szintillationsdetektor als Trigger des Messaufbaus als einfallende Teilchen gemessen werden und die eine Messung der Driftkammersignale auslösen, am häufigsten tatsächlich Elektronen und am seltensten Myonen oder andere zufällig einfallende Teilchen oder auch nur elektronisches Rauschen. Im Sinne des Experiments ist für dieses maximale Verhältnis dann also das Signal-zu-Rauschen-Verhältnis maximal und damit $U_\text{Szint}$ ideal.

Die für verschiedene Werte von $U_\text{Szint}$ gemessenen Daten sind in Tabelle \ref{tab:ratios_data} zu finden, die bestimmten Raten und Verhältnisse sind in Tabelle \ref{tab:ratios} aufgeführt. Hierbei ist zu beachten, dass die Raten ohne Präperat $n_\text{ohne}=n_\mu$ den Messungen in Abschnitt \ref{sec:myonenratios} entsprechen.

Das höchste Verhältnis $A=\num{11.70(50)}$ ergab sich bei der Einstellung $U_\text{Szint}=\SI{-1.750(18)}{\kilo\V}$, für die Einstellung der Betriebsparameter wurde daher $U_\text{Szint}$ auf diesen Wert eingestellt.

\begin{table}[h]
    \centering
    \begin{tabular}{|c||c|c||c|c|} \hline
    \input{../figs/ratios.table}
    \end{tabular}
    \caption{Messdaten der Spannungseinstellung des Szintillationsdetektors.}
    \label{tab:ratios_data}
\end{table}

\begin{table}[h]
    \centering
    \begin{tabular}{|c||c|c|c|} \hline
    \input{../figs/ratios_calc.table}
    \end{tabular}
    \caption{Berechnete Raten und Verhältnisse für die Spannungseinstellung des Szintillationsdetektors.}
    \label{tab:ratios}
\end{table}
